\documentclass[12pt]{article}
\usepackage[utf8]{inputenc}
\usepackage{microtype}
\usepackage{amsmath}
\usepackage{siunitx}
\usepackage[a4paper]{geometry}
\usepackage{hyperref}
\author{Thomas Orgis}
\title{Metronome acceleration}

\newcommand\de[0]{\text d}
\newcommand\ide[0]{\text d\,}
\newcommand\dd[2]{\frac{\de #1}{\de #2}}
\newcommand\ddq[2]{\frac{\de^2 #1}{\de {#2}^2}}
\newcommand\dnd[1]{\frac{\de}{\de #1}}
\newcommand\gkl[1]{\left( #1 \right)}

\begin{document}
\maketitle

\section*{Spacetime Analogy}

For some reason I keep confusing myself with the units in beat space and
the definition of an acceleration in reciprocal time only. How's that
with the usual motion over distance $s$ during time $t$?

The speed is
\begin{equation}
	v = \dd st
\;,
\end{equation}
while the acceleration is given as speed increase per covered distance,
\begin{equation}
	a = \dd vs = \dnd s \dd st
\;.
\end{equation}
This definition allows to compute the speed after a certain covered
distance,
\begin{equation}
\begin{aligned}
\label{eq-speed}
	v(s) &= v_0 + \int\limits_{s_0}^s a(s') \de s'
\\
\stackrel{a = \text{const}}{\Leftrightarrow}
	v(s) &= v_0 + a(s-s_0)
\;.
\end{aligned}
\end{equation}

That can be again regarded as a differential equation linking
space and time increments, to be solved for the time difference
corresponding to two points in space,
\begin{equation}
\begin{aligned}
	v =
	\dd st &= v_0 + a(s-s_0)
\\
\Leftrightarrow
	\de t &= \frac 1{v_0 + a(s-s_0)} \de s
\\
\Leftrightarrow
	t_1 - t_0 &=
	\int\limits_{s_0}^{s_1}  \frac 1{v_0 + a(s-s_0)} \de s
\\\text{\ldots substituting \autoref{eq-speed} \ldots}
\\
	v = v_0 + a(s-s_0)
	&\Rightarrow \dd vs = a
	\Rightarrow \de s = \frac 1a \de v
\\
\Leftrightarrow
        t_1 - t_0 &=                           
        \frac 1a \int\limits_{v(s_0)}^{v(s_1)}  \frac 1v \de s
\;.
\end{aligned}
\end{equation}
Hence, the duration of a certain distance interval in time is given by
\begin{equation}
\label{eq-deltat}
        \Delta t = t_1 - t_0
	= \frac 1a \ln \frac {v(s_1)}{v(s_0)}
\;.
\end{equation}

How should one interpret the dimension of the acceleration?
The unit \si{s^{-1}} could be seen as \si{(m/s) / m},
\si{h^{-1}} as  \si{(km/h) / km}. One mixed variant
could be \si{(km/h) / m}, which boils down to \SI{1000}{h^{-1}}.


\section*{Real Deal}

The analogy can be carried forward, identifying the distance $s$ with
the position in the score, counted in bars, beats, whatever you like.
A sensible unit might be a whole note, denoted with the unit symbol \si{n}.
A quarter note equals $\Delta s_{\text note} = \SI{1/4}{n}$, an eigth $\SI{1/8}{n}$.
A bar in 6/8 meter has length $\Delta s_{\text{bar}} = \SI{3/4}{n}$,
while in 4/4 meter it's an unsurprising $\Delta s_{\text{bar}} = \SI{1}{n}$.

Dominic explained the meaning of the acceleration parameter $c$ as such:
After $c$ bars, the tempo increased by \SI{1}{bpm}.
The tempo is always related to quarter notes, so that one bar of three
quarters length at constant tempo $v=\SI{120}{bpm}$, be it either in
$3/4$ or in $6/8$ meter, takes \SI{1.5}s to complete. With the bar length

The acceleration value, based on quarter beats per minute, is
given by
\begin{equation}
\label{eq-acc1}
	a \cdot (c \Delta s_{\text{bar}})
	=
	\SI{1}{bpm}
	=
	1 \frac {\SI{1/4}n}{\si{min}}
\Leftrightarrow
	\frac a { \si{min^{-1}} } =
	\frac 1c \frac { \SI{1/4}n }{ \Delta s_{\text{bar}} }
\;.
\end{equation}
The tempo $v(s)$ at a certain point into the score is given by
\autoref{eq-speed} using this acceleration. Also, a time interval
for a certain beat can be computed using \autoref{eq-deltat}
and the tempo values at the beginning of the converned beat ($s_0$)
and the beginning of the next one ($s_1$).

It is evident from \autoref{eq-deltat} that all this is only applicable
if there indeed is a non-zero acceleration $a$ and the tempo is
always non-zero.
The latter point is illustrative in the spacetime analogy with the
conventional definition of a constant acceleration per time interval,
\begin{align}
	a &= \dd vt = \ddq st
\;,
\end{align}
where everything is defined nicely. It {\em is} a difference whether
the acceleration is defined using elapsed time or covered
distance.
If you start out with $v_0 = 0$ (in whatever units), an acceleration
that is defined to increase $v$ with each increment of distance cannot
take hold, as there is no infinitisemal distance to work on at the
beginning. You can't get going if you're not going.
The solution makes sense: In that case, there is no solution.

Negative acceleration is of course allowed, as long as it is
non-zero. For constant tempo, the trivial
\begin{equation}
	\Delta t = \frac{\Delta s}v
\end{equation}
applies.

\subsection*{Derived Acceleration}

Actually, the usual notion inside a score is to change tempo from $v_0$ to
$v_1$ within a certain number of bars. In place of \autoref{eq-acc1}, the
acceleration then simply is given by the ratio of the tempo difference and
the corresponding distance in the score,
\begin{equation}
\label{eq-acc2}
	a = \frac {v(s_1)-v(s_0)}{s_1 - s_0}
	= \frac{\Delta v}{\Delta s}
\;.
\end{equation}

This the obvious sensible default to use when a part length in bars is given
without acceleration.

\subsection*{Mean Tempo}

To compare with Dominic's empirical mean tempo for a beat,
\begin{equation}
\bar v = \frac{v_0 - v_1}{\ln \frac{v_0}{v_1}}
= -\frac{ v_1-v_0 }{ \ln \frac{v_0}{v_1} }
=  \frac{ v_1-v_0 }{ \ln \frac{v_1}{v_0} }
\end{equation}
(written to avoid direct logarithms of quantities with units attached),
the acceleration in \autoref{eq-deltat} can be replaced by
\autoref{eq-acc2}. This yields
\begin{equation}
	\Delta t = \frac{\Delta s}{\Delta v}
	\ln \frac {v_1}{v_0}
\Leftrightarrow
\bar v = \frac{\Delta s}{\Delta t}
= \frac{\Delta v}{\ln \frac {v_1}{v_0}}
\;,
\end{equation}
which is the same as Dominic's expression.

\section*{In Terms of Time}

% Note: I'm writing this sitting in a train called "Metronom"!
One could also define the metronome tempo acceleration in terms of time.
Instead of being given in units of inverse time, it would be
beats-per-minute per minute.
But how easy is it to solve for the length of a beat $\Delta s$ then?

\begin{equation}
\begin{aligned}
v(t) &= v_0 + a(t - t_0)
\\
\Delta s &= \int\limits_{t_0}^{t_1} v_0 + a(t - t_0) \de t
\\
&= v_0 \Delta t + \frac a2 (\Delta t)^2
\\
0 &= (\Delta t)^2 + \frac 2a v_0 \Delta t - \frac 2a \Delta s
\end{aligned}
\end{equation}
\begin{equation}
\Delta t = - \frac {v_0}a \pm \sqrt{ \gkl{\frac{v_0}a}^2 + \frac 2a \Delta s }
\end{equation}
This is straightfoward, allowing for the choice of the positive $\Delta t$.
Now, what is desired? Constant acceleration in space (beat count) or in time?

\end{document}
